\title{The Era of 1-bit LLMs: All Large Language Models are in 1.58 Bits}
Memory usage shows a parallel trend, as the fixed full-precision embedding constitutes a diminishing fraction of total memory for more extensive models. To analyze resource efficiency, we benchmarked GPU memory usage and inference latency for both \text{LLaMA LLM} and \text{BitNet b1.58}. In line with perplexity trends, end-task outcomes further demonstrate that \text{BitNet b1.58} 3.9B surpasses \text{LLaMA LLM} 3B, maintaining reduced computational cost in memory and latency. \section{BitNet b1.58}

\text{BitNet b1.58} builds upon the BitNet framework, a variant of Transformer architecture in which standard \emph{nn.Linear} modules are substituted with \emph{BitLinear} counterparts. The trajectory of research has been toward progressively reduced bit-width, such as moving to 4-bit models~\cite{gptq, awq}. Both latency and memory metrics were collected using a 2-bit kernel, implying potential for further cost reduction through additional optimization. \text{BitNet b1.58} inherits all essential advantages from its 1-bit predecessor: it employs an innovative computation strategy that minimizes reliance on multiplication in matrix operations, thereby enabling aggressive optimization, and its energy usage is on par with the original, while it further reduces memory requirements, increases throughput, and lowers latency relative to FP16 LLMs. \textbf{\text{BitNet b1.58} introduces a novel scaling law concerning model performance relative to inference efficiency.} The outcomes shown in Figures~\ref{fig:latency-memory-energy} and~\ref{fig:energy} enable us to draw parallels between 1.58-bit and 16-bit model sizes in terms of practical equivalence:

\begin{itemize}
    \item The 13B BitNet b1.58 model surpasses the 3B FP16 LLM in latency, memory consumption, and energy efficiency. These metrics were obtained with the FasterTransformer\footnote{\url{https://github.com/NVIDIA/FasterTransformer}} framework, which is highly optimized for large language model inference on GPUs. Besides computation, transferring parameters from DRAM to on-chip accelerator memory (e.g., SRAM) also contributes significantly to inference cost. A promising strategy to mitigate these limitations involves applying post-training quantization techniques to construct lower-bit inference models~\cite{smoothquant, gptq, quip, quip_sharp}. \textbf{New Hardware for 1-bit LLMs}

Recent developments, such as Groq\footnote{\url{https://groq.com/}}, highlight significant strides in developing hardware (for instance, LPUs) tailored for LLM workloads. For benchmarks providing both accuracy and normalized accuracy, we averaged these metrics. To evaluate how well \text{BitNet b1.58} scales with respect to training data, we trained it on 2T tokens, utilizing the data pipeline adopted for StableLM-3B~\cite{StableLM-3B-4E1T}, a leading open-source model at the 3B scale. Approaches that attempt to increase SRAM size to boost throughput generally face much higher costs than DRAM. Moreover, \text{BitNet b1.58} offers two significant improvements. Notably, from the 3B scale onwards, \text{BitNet b1.58} achieves performance on par with its full-precision counterpart. Figure~\ref{fig:energy} provides a breakdown of the energy usage, revealing that the computation in \text{BitNet b1.58} predominantly involves INT8 addition, in contrast to \text{LLaMA LLM}, which utilizes both FP16 addition and FP16 multiplication. These results indicate that \text{BitNet b1.58} represents a Pareto improvement relative to current state-of-the-art large language models. \paragraph{LLaMA-alike Components.}

The design of LLaMA~\cite{llama, llama2} has become the prevalent standard for open-source large language models. Furthermore, the compatibility of 1.58-bit LLMs with CPUs—ubiquitous in edge and mobile environments—means that \text{BitNet b1.58} can run efficiently, delivering further improvements in usability and performance. First, it achieves enhanced expressive power through explicit feature filtering enabled by the additional zero value in the parameter set, which contributes to superior performance relative to conventional 1-bit LLMs. Evaluations adhered to the \emph{lm-evaluation-harness}\footnote{\url{https://github.com/EleutherAI/lm-evaluation-harness}} standard pipeline. We assessed both models on a composite benchmark comprising Winogrande~\cite{winoGrande}, PIQA~\cite{piqa}, SciQ~\cite{sciq}, LAMBADA~\cite{lambada}, and ARC-easy~\cite{arc}, reporting zero-shot accuracy in Table~\ref{tab:2t}. Ideally, consolidating entire models onto a single chip would eliminate this overhead altogether. StableLM 3b results at the 2T scale were sourced from its official technical report. Notably, \text{BitNet b1.58} 70B outperforms the \text{LLaMA LLM} baseline by achieving a 4.1-fold increase in speed. Despite a notable decrease in computational FLOPs, these models are often constrained by substantial memory usage and significant communication costs between chips, which hamper broader deployment. There is potential for further reduction—compressing activations to 4 bits or below for 1.58-bit LLMs—suggesting a direction for subsequent research. \end{document} As reported in Table~\ref{tab:throughput}, the 70B \text{BitNet b1.58} model accommodated up to 11-fold larger batch sizes compared to \text{LLaMA LLM}, which translated to an 8.9-fold increase in throughput. This model is trained from the beginning using weights quantized to 1.58 bits and activations at 8 bits. Instead, per token, activations are rescaled within $[-Q_b, Q_b]$, thereby eliminating zero-point quantization. The introduction of 1.58-bit LLMs offers solutions to these barriers. In particular, the design incorporates RMSNorm~\cite{rmsnorm}, SwiGLU~\cite{swiglu}, rotary embedding~\cite{rope}, and eliminates all bias parameters. Memory usage, particularly from KV caches during inference, poses a primary obstacle for long sequence handling. This adjustment streamlines both code implementation and hardware optimization with a negligible impact on empirical performance. \paragraph{Quantization Function.} To enforce weight values to be one of -1, 0, or +1, an \emph{absmean} quantization strategy is used. Validation perplexity metrics were collected on WikiText2~\cite{wikitext2} and C4~\cite{c4}. Table~\ref{tab:ppl} provides an overview of both perplexity scores and computational costs associated with \text{BitNet b1.58} and \text{LLaMA LLM}. By lowering the precision of model parameters and activations, the memory footprint and computational load of LLMs can be drastically decreased. The findings indicate that as model size increases, the relative energy efficiency advantage of \text{BitNet b1.58} over the FP16 \text{LLaMA LLM} grows. Given that power often serves as a limiting factor in the computational capabilities of modern hardware, this substantial reduction in energy expenditure can be directly leveraged for faster model inference. \item The 70B BitNet b1.58 model achieves greater efficiency in latency, memory, and energy than the 13B FP16 LLM. Furthermore, at the 3.9B model scale, \text{BitNet b1.58} delivers a 2.4$\times$ increase in speed, requires 3.32$\times$ less memory, and yields better performance compared to the \text{LLaMA LLM} 3B model. The findings indicate that at the 3B scale, \text{BitNet b1.58} reaches parity with the full-precision \text{LLaMA LLM} in perplexity, while offering a 2.71$\times$ speedup and reducing GPU memory usage by a factor of 3.55. Building on these advancements, we propose future efforts be directed toward creating specialized hardware and systems fine-tuned for the unique computational requirements of 1-bit LLMs, leveraging the novel computation paradigms introduced by BitNet~\cite{bitnet}. Recent advancements in 1-bit model architectures, exemplified by BitNet~\cite{bitnet}, indicate a promising strategy for decreasing the computational demands of LLMs without compromising accuracy. By incorporating a 0 into the original 1-bit BitNet parameterization, the effective precision becomes 1.58 bits within the binary framework. Traditional LLMs typically utilize 16-bit floating-point representations (such as FP16 or BF16), and their main computational bottleneck arises from extensive matrix multiplications involving floating-point additions and multiplications. \end{equation}

The approach to activation quantization is kept consistent with BitNet's method; however, we no longer normalize activations prior to applying non-linearities to fit into the interval $[0, Q_b]$. Further, the overall energy required for processing 512 tokens was documented. \textbf{LLMs on Edge and Mobile}

Implementing 1.58-bit LLMs opens up enhanced possibilities for deploying language models on edge and mobile hardware. This improvement is attributed to the escalating computational time required for \emph{nn.Linear} operations with larger models. The data suggest that the accuracy gap between \text{BitNet b1.58} and \text{LLaMA LLM} diminishes with growing model size. Table~\ref{tab:zero-shot} details zero-shot accuracy results on downstream tasks. In this study, we present a novel 1-bit LLM variant called \textbf{\text{BitNet b1.58}}, wherein each model parameter is ternary, allowed to take on values from the set \{-1, 0, 1\}. \paragraph{Memory and Latency}

The model was extended to larger configurations, specifically 7B, 13B, and 70B parameters, and the associated costs were assessed. This procedure involves normalizing the weight matrix by dividing by its mean absolute value, followed by rounding each entry to the closest integer from \{-1, 0, +1\}:

\begin{equation}
    \widetilde{W} = \mathrm{RoundClip}(\frac{W}{\gamma + \epsilon}, -1, 1),
\end{equation}

\begin{equation}
    \mathrm{RoundClip}(x, a, b) = \max(a, \min(b, \mathrm{round}(x))),
\end{equation}

\begin{equation}
    \gamma = \frac{1}{nm}\sum_{ij} |W_{ij}|. \section{Discussion and Future Work}

\textbf{1-bit Mixture-of-Experts (MoE) LLMs}

Mixture-of-Experts (MoE) architectures are recognized for their efficiency in scaling Large Language Models (LLMs). For an equitable assessment, both models underwent pre-training on the RedPajama dataset~\cite{redpajama}, utilizing 100 billion tokens. These choices enable straightforward adoption of \text{BitNet b1.58} into major open-source frameworks, such as Huggingface, vLLM~\cite{vllm}, and llama.cpp\footnote{\url{https://github.com/ggerganov/llama.cpp}}, requiring only minor modifications. \text{BitNet b1.58} advances toward inherent long sequence support by shrinking activation precision from 16 bits to 8 bits, which enables models to double their context length within the same hardware limitations. \paragraph{Throughput}

To assess throughput, we benchmarked \text{BitNet b1.58} against \text{LLaMA LLM} with 70B parameters, deploying both models on a pair of 80GB A100 GPUs. This approach achieves parity with full-precision (e.g., FP16 or BF16) Transformer LLMs—matched by both perplexity and performance on downstream tasks—while providing substantial gains in latency, memory usage, throughput, and energy efficiency. 
\begin{abstract}
Recent advancements, such as BitNet~\cite{bitnet}, are ushering in a transformative era characterized by 1-bit Large Language Models (LLMs). \section{Results}

We conducted a comparison between \text{BitNet b1.58} and our reproduced FP16 \text{LLaMA LLM} across multiple parameter scales. In addition, it establishes a foundation for alternative computational paradigms and paves the way for dedicated hardware architectures tailored specifically to 1-bit LLM workloads. \item The 30B BitNet b1.58 model outperforms the 7B FP16 LLM in terms of latency, memory requirements, and energy usage. This efficiency not only curtails loading time and associated expenses when moving weights from DRAM but also results in more expedient and cost-effective inference. To further engage with the open-source ecosystem, \text{BitNet b1.58} is constructed utilizing LLaMA-alike architectural features. We employed pipeline parallelism~\cite{gpipe} to enable execution of the 70B parameter \text{LLaMA LLM}. The compact memory footprint permits deployment with fewer hardware units, thus reducing the device count for hosting MoE models. \paragraph{Energy}

An analysis of arithmetic energy consumption was conducted for both \text{BitNet b1.58} and \text{LLaMA LLM}, centering primarily on matrix multiplication operations, given their substantial impact on the computational budget of large language models. Relative to the original BitNet, several alterations have been made, which are outlined below. With the diminished memory and energy demands of 1.58-bit LLMs, these models become feasible on resource-constrained platforms, facilitating novel use cases previously unattainable. For \text{BitNet b1.58}, integration with the 2-bit kernel from Ladder~\cite{ladder} was also included. Additionally, lower memory requirements translate to reduced inter-device activation transfer, substantially cutting network overhead. 1-bit LLMs, in comparison to their full-precision counterparts, demand considerably less memory both in terms of required capacity and bandwidth. \textbf{Native Support of Long Sequence in LLMs}

The growing significance of LLMs brings an increased need for processing extended sequences. \end{abstract}

\section{The Era of 1-bit LLMs}

The field of AI has witnessed an explosive increase in the scope and power of Large Language Models (LLMs) in recent years. This is largely due to the rising contribution of \emph{nn.Linear} layers to overall computation in larger models, whereas the proportional impact of other modules decreases. These systems excel across diverse natural language processing applications; nevertheless, the escalating model sizes complicate real-world deployment and heighten concerns related to ecological and financial costs owing to intensive energy requirements. Notably, the 1.58-bit LLM framework offers a novel scaling law and methodology for training the next wave of LLMs that prioritize both efficiency and accuracy. Consequently, the functionality of edge and mobile devices can be significantly expanded, ushering in innovative applications of LLMs. Second, our experimental findings demonstrate that, with comparable configurations (including model scale and amount of training data), \text{BitNet b1.58} can attain parity with full-precision (FP16) baselines in both perplexity and task-oriented evaluations, beginning with the 3B parameter regime. Drawing from the energy modeling framework in~\cite{energycost, pokebnn}, \text{BitNet b1.58} achieves a 71.4-fold reduction in energy consumption for matrix multiplication on 7nm hardware. As depicted in Figure~\ref{fig:latency-memory-energy}, both latency and memory usage exhibit systematic patterns as model size increases, with observed speed-ups becoming more pronounced for larger models. By progressively increasing the batch size up to the point where the GPU memory was saturated, while maintaining a fixed sequence length of 512, we observed substantial differences. Despite their popularity in industrial settings, post-training quantization methods still fall short of optimality. \end{itemize}

\paragraph{Training with 2T Tokens}

The volume of training tokens plays a pivotal role in shaping LLM capabilities. Our analysis demonstrates that \text{BitNet b1.58} outperforms across all evaluation tasks, highlighting the robust generalization abilities of LLMs using 1.58-bit quantization. By contrast, BitNet’s matrix multiplications are based solely on integer addition, which drastically reduces energy consumption. Time per output token was chosen as the principal measure of inference efficiency. The models' zero-shot capabilities were assessed over diverse language understanding benchmarks, such as ARC-Easy~\cite{arc}, ARC-Challenge~\cite{arc}, Hellaswag~\cite{hellaswag}, Winogrande~\cite{winoGrande}, PIQA~\cite{piqa}, OpenbookQA~\cite{openbookqa}, and BoolQ~\cite{boolq}. In this study, we present a new 1-bit LLM, denoted as \textbf{\text{BitNet b1.58}}, where each parameter (or weight) of the LLM adopts one of three values from the set \{-1, 0, 1\}. Typically, such devices are restricted by modest memory and computational capabilities, limiting feasible model size and performance.